\textbf{Phần 1. Trắc nghiệm lựa chọn}
\setcounterpageref{ex}{0}

\begin{ex} %[3D1V61]
Làm tròn số $47\,620$ đến hàng nghìn ta được số
\choice
{\True $48\,000$}
{$47\,000$}
{$50\,000$}
{$47\,600$}
\loigiai{
Chữ số hàng trăm là 6 ($6 \ge 5$) nên làm tròn lên hàng nghìn: $48\,000$.
\textbf{Đáp án của câu hỏi là A}
}
\end{ex}

\begin{ex} %[3D1V61]
Làm tròn số $73\,480$ đến hàng chục nghìn ta được số
\choice
{$80\,000$}
{$73\,000$}
{\True $70\,000$}
{$74\,000$}
\loigiai{
Chữ số hàng nghìn là 3 ($3 < 5$) nên làm tròn xuống: $70\,000$.
\textbf{Đáp án của câu hỏi là C}
}
\end{ex}

\begin{ex} %[3D1V61]
Số bé nhất có năm chữ số mà khi làm tròn đến hàng nghìn được $65\,000$ là
\choice
{$65\,499$}
{\True $64\,500$}
{$64\,000$}
{$64\,550$}
\loigiai{
Các số làm tròn đến hàng nghìn thành 65 000 là từ 64 500 đến 65 499. Số bé nhất là $64\,500$.
\textbf{Đáp án của câu hỏi là B}
}
\end{ex}

\begin{ex} %[3D1V61]
Số lớn nhất có năm chữ số mà khi làm tròn đến hàng chục nghìn được $80\,000$ là
\choice
{$85\,000$}
{$79\,999$}
{$84\,500$}
{\True $84\,999$}
\loigiai{
Các số làm tròn đến hàng chục nghìn thành 80 000 là từ 75 000 đến 84 999. Số lớn nhất là $84\,999$.
\textbf{Đáp án của câu hỏi là D}
}
\end{ex}

\begin{ex} %[3D1V61]
Dân số một huyện là $96\,720$ người. Báo cáo làm tròn dân số đến hàng chục nghìn là khoảng
\choice
{\True $100\,000$ người}
{$90\,000$ người}
{$97\,000$ người}
{$96\,000$ người}
\loigiai{
Chữ số hàng nghìn là 6 ($6 \ge 5$) nên làm tròn lên: $100\,000$ người.
\textbf{Đáp án của câu hỏi là A}
}
\end{ex}

\begin{ex} %[3D1V61]
Làm tròn số $54\,890$ đến hàng nghìn được số $A$, làm tròn đến hàng chục nghìn được số $B$. Giá trị của $B - A$ là
\choice
{$4000$}
{$5000$}
{\True $-5000$ (hoặc hiệu chênh lệch là $5000$)}
{$6000$}
\loigiai{
$A = 55\,000$; $B = 50\,000$. Chênh lệch giữa hai số là: $55\,000 - 50\,000 = 5000$.
\textbf{Đáp án của câu hỏi là C}
}
\end{ex}

\begin{ex} %[3D1V61]
Số nào sau đây khi làm tròn đến hàng nghìn thì được số $32\,000$?
\choice
{$32\,600$}
{\True $32\,450$}
{$31\,400$}
{$32\,550$}
\loigiai{
Số 32 450 có hàng trăm là 4 ($4 < 5$) nên làm tròn thành $32\,000$.
\textbf{Đáp án của câu hỏi là B}
}
\end{ex}

\begin{ex} %[3D1V61]
Khi làm tròn số $95\,000$ đến hàng chục nghìn ta được số nào?
\choice
{$90\,000$}
{$95\,000$}
{$96\,000$}
{\True $100\,000$}
\loigiai{
Chữ số hàng nghìn là 5 ($5 \ge 5$) nên làm tròn lên $100\,000$.
\textbf{Đáp án của câu hỏi là D}
}
\end{ex}

\begin{ex} %[3D1V61]
Quy tắc làm tròn đến hàng nghìn: ta quan sát chữ số ở hàng nào?
\choice
{\True Hàng trăm}
{Hàng chục}
{Hàng đơn vị}
{Hàng chục nghìn}
\loigiai{
Để làm tròn đến hàng nghìn, ta xét chữ số ở hàng trăm liền kề bên phải.
\textbf{Đáp án của câu hỏi là A}
}
\end{ex}

\begin{ex} %[3D1V61]
Một công ty sản xuất được $68\,500$ sản phẩm. Làm tròn đến hàng nghìn ta được số sản phẩm là
\choice
{$68\,000$ sản phẩm}
{$70\,000$ sản phẩm}
{\True $69\,000$ sản phẩm}
{$68\,600$ sản phẩm}
\loigiai{
Hàng trăm là 5 nên làm tròn lên $69\,000$ sản phẩm.
\textbf{Đáp án của câu hỏi là C}
}
\end{ex}

\begin{ex} %[3D1V61]
Có bao nhiêu số có năm chữ số khi làm tròn đến hàng nghìn thì được số $40\,000$?
\choice
{$999$ số}
{\True $1000$ số}
{$500$ số}
{$100$ số}
\loigiai{
Các số từ 39 500 đến 40 499 gồm: $40\,499 - 39\,500 + 1 = 1000$ số.
\textbf{Đáp án của câu hỏi là B}
}
\end{ex}

\begin{ex} %[3D1V61]
Khoảng cách từ Trái Đất đến một vệ tinh là $35\,820\text{ km}$. Làm tròn đến hàng chục nghìn, khoảng cách đó là
\choice
{$30\,000\text{ km}$}
{$36\,000\text{ km}$}
{$35\,000\text{ km}$}
{\True $40\,000\text{ km}$}
\loigiai{
Hàng nghìn là 5 nên làm tròn lên $40\,000\text{ km}$.
\textbf{Đáp án của câu hỏi là D}
}
\end{ex}

\textbf{Phần 2. Trắc nghiệm đúng sai}
\setcounterpageref{ex}{0}

\begin{ex} %[3D1V61]
Xét tính đúng sai của việc làm tròn số đến hàng nghìn và hàng chục nghìn:
\choiceTF
{\True Khi làm tròn đến hàng chục nghìn, nếu chữ số hàng nghìn là 5 thì ta làm tròn lên.}
{Số 44 999 khi làm tròn đến hàng chục nghìn thì được 50 000.}
{\True Số 85 000 làm tròn đến hàng chục nghìn được 90 000.}
{Làm tròn số đến hàng nghìn luôn cho kết quả lớn hơn số ban đầu.}
\loigiai{
\begin{itemize}
    \item Ý a) Đúng: Quy tắc chuẩn.
    \item Ý b) Sai: Sai vì chữ số hàng nghìn là 4 nên làm tròn thành 40 000.
    \item Ý c) Đúng: Vì hàng nghìn là 5.
    \item Ý d) Sai: Sai vì nếu làm tròn xuống thì kết quả nhỏ hơn.
\end{itemize}
\begin{itemize}
    \item a) \textbf{Đúng}.
    \item b) \textbf{Sai}.
    \item c) \textbf{Đúng}.
    \item d) \textbf{Sai}.
\end{itemize}
}
\end{ex}

\begin{ex} %[3D1V61]
Cho số $67\,485$:
\choiceTF
{\True Làm tròn đến hàng nghìn được số 67 000.}
{\True Làm tròn đến hàng chục nghìn được số 70 000.}
{Làm tròn đến hàng trăm được số 67 400.}
{\True Hiệu giữa số làm tròn đến hàng chục nghìn và hàng nghìn là: $70\,000 - 67\,000 = 3000$.}
\loigiai{
\begin{itemize}
    \item Ý a) Đúng: Hàng trăm là 4 làm tròn xuống.
    \item Ý b) Đúng: Hàng nghìn là 7 làm tròn lên.
    \item Ý c) Sai: Sai vì hàng chục là 8 nên làm tròn thành 67 500.
    \item Ý d) Đúng: Phép tính đúng.
\end{itemize}
\begin{itemize}
    \item a) \textbf{Đúng}.
    \item b) \textbf{Đúng}.
    \item c) \textbf{Sai}.
    \item d) \textbf{Đúng}.
\end{itemize}
}
\end{ex}

\textbf{Phần 3. Trả lời ngắn}
\setcounterpageref{ex}{0}

\begin{ex} %[3D1V61]
Làm tròn số 58 640 đến hàng nghìn.
\shortans{59000}
\loigiai{
\textbf{Đáp số:} $59000$\\
Hàng trăm là 6 nên làm tròn thành 59 000.
}
\end{ex}

\begin{ex} %[3D1V61]
Làm tròn số 82 500 đến hàng chục nghìn.
\shortans{80000}
\loigiai{
\textbf{Đáp số:} $80000$\\
Hàng nghìn là 2 nên làm tròn thành 80 000.
}
\end{ex}

\begin{ex} %[3D1V61]
Tìm số tự nhiên lớn nhất có 5 chữ số khi làm tròn đến hàng nghìn được 74 000.
\shortans{74499}
\loigiai{
\textbf{Đáp số:} $74499$\\
Số lớn nhất là 74 499.
}
\end{ex}

\begin{ex} %[3D1V61]
Tìm số tự nhiên bé nhất có 5 chữ số khi làm tròn đến hàng chục nghìn được 50 000.
\shortans{45000}
\loigiai{
\textbf{Đáp số:} $45000$\\
Số bé nhất là 45 000.
}
\end{ex}

\begin{ex} %[3D1V61]
Làm tròn số 99 600 đến hàng nghìn.
\shortans{100000}
\loigiai{
\textbf{Đáp số:} $100000$\\
Hàng trăm là 6 làm tròn lên: $99\,000 + 1000 = 100\,000$.
}
\end{ex}

\begin{ex} %[3D1V61]
Một sân vận động có 43 780 khán giả. Làm tròn đến hàng nghìn, số khán giả là bao nhiêu?
\shortans{44000}
\loigiai{
\textbf{Đáp số:} $44000$\\
Hàng trăm là 7 nên làm tròn thành 44 000.
}
\end{ex}
